%======================================================================
\section{Cómo entra el riesgo en cada formulación}\label{sec:riesgo-como-entra}
%======================================================================
Esta sección escribe cada peldaño de la escalera como una modificación explícita del QP de N2,
dice sobre qué se monta, cuánto cuesta y qué supone, y explica por qué la capa se puso sobre N2 y
no sobre N3.

\subsection{El núcleo: N2 en lazo cerrado}
En cada consulta, N2 resuelve un QP sobre el modelo linealizado a lo largo de la trayectoria
prevista (sección~\ref{sec:niveles}). En forma condensada, $\hat y=\hat y^0+Gu$, con $G$ la
sensibilidad de la PCT predicha al caudal (triangular inferior, $G_{ki}<0$: más caudal enfría):
\begin{align*}
\text{(N2)}\qquad\min_{u,\,r}\quad & \underbrace{\|\hat y^0+Gu+d\mathbf 1-y_{\rm ref}\mathbf 1\|^2}_{\text{seguimiento}}
+\underbrace{R\,\|\Delta u\|^2}_{\text{suavidad}}+\underbrace{\lambda_{\rm pct}\|r\|^2}_{\text{límite }U_y}\\
\text{s.a.}\quad & 0\le u_k\le\mdot_{\max},\qquad |u_k-u_{k-1}|\le\Delta\mdot,\qquad
\hat y^0_k+(Gu)_k+d-U_y\le r_k,\quad r_k\ge0,
\end{align*}
con $d$ el sesgo medido, $U_y=1244$~K, $\mdot_{\max}=40$~kg/s, $\Delta\mdot=5$~kg/s por paso y
$R=0.1$. Es el peldaño 0. Ningún término contiene la falla: la predicción $\hat y^0$ supone que
el ACS entrega lo que se le pide.

\begin{figure}[H]\centering
\resizebox{\linewidth}{!}{%
\begin{tikzpicture}[
  core/.style={draw=accent, line width=0.9pt, rounded corners=3pt, fill=soft, align=center,
               font=\small, minimum height=22mm, text width=38mm},
  plug/.style={draw=accent, line width=0.5pt, rounded corners=2pt, fill=white, align=center,
               font=\footnotesize, minimum height=11mm, text width=37mm},
  arr/.style={-{Stealth[length=2mm]}, line width=0.6pt, draw=ink},
  lab/.style={font=\scriptsize, text=black!65, align=center, fill=white, inner sep=1pt}]
\node[core] (n2) at (0,0) {\textbf{QP de N2}\\seguimiento + $R\|\Delta u\|^2$\\caja y rampa duras\\(peldaño 0)};
\node[plug] (dsm) at (-6.2,1.6) {\textbf{DSM} (peldaño 1)\\término $\lambda[\mathrm{DSM}_k-\tau]_+^2$};
\node[plug] (cota) at (-6.2,-1.6) {\textbf{Techo} $1220-m$ (peldaño 2)\\$m$: cuantil, CVaR, Wilks};
\node[plug] (saa) at (6.2,1.6) {\textbf{Escenarios} (peldaño 2)\\$\rho\,\mathbb E[(\text{pico}-1220)_+]$, SAA};
\node[plug] (esc) at (6.2,-1.6) {\textbf{Un escenario} (peldaño 3)\\$\hat\delta=\Phi(\theta,x)$ aprendido};
\draw[arr] (dsm) -- node[lab, above, sloped]{costo} (n2);
\draw[arr] (cota) -- node[lab, below, sloped]{26 restricciones} (n2);
\draw[arr] (saa) -- node[lab, above, sloped]{$26\times40$ holguras} (n2);
\draw[arr] (esc) -- node[lab, below, sloped]{26 holguras} (n2);
\node[lab] (rom) at (-6.2,-3.4) {ROM sin ACS $\Rightarrow\Delta T(t)$};
\draw[arr, dashed] (rom) -- (cota);
\node[lab] (dist) at (6.2,-3.4) {datos de lazos cerrados $\Rightarrow\theta$};
\draw[arr, dashed] (dist) -- (esc);
\end{tikzpicture}}
\caption{Las capas de riesgo se enchufan al mismo QP. Todas siguen siendo un QP convexo del
tamaño de N2 o un poco mayor; ninguna cambia el modelo de predicción.}
\label{fig:riesgo-enchufes}
\end{figure}

\subsection{Las formulaciones, una por una}
\begin{table}[H]\centering\scriptsize\setlength{\tabcolsep}{2.5pt}
\begin{tabularx}{\linewidth}{@{}>{\raggedright\arraybackslash}p{24mm} >{\raggedright\arraybackslash}p{50mm} >{\raggedright\arraybackslash}p{16mm} Y >{\raggedright\arraybackslash}p{14mm}@{}}\toprule
Formulación & Qué se agrega a N2 & Base & Tamaño y tiempo (p99) & Fase\\\midrule
\multicolumn{5}{@{}l}{\textit{Peldaño 0: determinista}}\\
N1 / N2 / N3 & nada (límite $U_y$ blando, inactivo) & --- & 77 var.; 0.03--0.05~s (N3: NLP de 519 var., 0.2~s) & F2--F6\\
\multicolumn{5}{@{}l}{\textit{Peldaño 1: margen determinista}}\\
DSM de la tesis & $+\lambda_{\rm dsm}\sum_k[\mathrm{DSM}_k-\tau]_+^2$, $\mathrm{DSM}_k=\frac{\hat y_k+d-y_{\rm ref}}{U-y_{\rm ref}}$ & N2 & +26 holguras; igual & F5b.1b\\
\multicolumn{5}{@{}l}{\textit{Peldaño 2: métricas probabilísticas}}\\
Pronóstico gaussiano & $\hat y_k+z_{1-\varepsilon}\sigma_k\le U$ & N1 & igual & F4 (descartada)\\
SAA de pronóstico & $\frac1M\sum_jJ(u_1,u^{(j)}_{2..K};\hat d^{(j)})$, $M=59$ & N3 & 30\,000 var.; 2~s & F4\\
Cota con cuantil & $\hat y_k+d\le1220-\Delta T(q_{1-\alpha})+s_k$, $+\lambda_r\sum s_k$ & N2 & +26 holguras, 1 simulación ROM; 0.1~s & F5b.1c\\
CVaR$_\beta$ & ídem con $m=\mathrm{CVaR}_\beta[\Delta T(D)]$ (200 cuantiles) & N2 & ídem & F5b.1c, F6\\
Wilks 59 & ídem con $m=\Delta T(\max_{s\le59}D_s)$, sorteo nuevo por consulta & N2 & ídem & F5b.1c, F6\\
Pérdida esperada & $+\frac{\rho}{40}\sum_{k,s}v_{ks}$, $v_{ks}\ge\hat y_k+d+\Delta T(D_s)-1220$, $v\ge0$ & N2 & +1040 var.; 0.3--0.4~s & F5b.1c, F6, F9\\
Fallas parciales & escenarios $(j,\eta,D)$: $\hat y^{(s)}_k=\hat y^0_k+\sum_{i<j}G_{ki}u_i+\eta\sum_{i\ge j}G_{ki}u_i$ & N2 & +26 por escenario & F5b.1d\\
Cota con $\eta$ & cuantil con $\sum_\eta p_\eta P(\Delta T_\eta(D)>m)=\alpha$ & N2 & 4 curvas ROM & F5b.1d\\
\multicolumn{5}{@{}l}{\textit{Peldaño 3: escenario único}}\\
Escenario del teorema & segunda etapa con $\xi^*$ (un margen por paso) & N2 & 26 holguras; requiere el SAA antes & F9 (E2)\\
Escenario aprendido & $+\rho\sum_k(\hat y_k+d+\hat\delta-1220)_+$, $\hat\delta=\Phi(\theta,x)$ & N2 & 26 holguras; 0.03--0.07~s & F9 (E4--E5)\\
\multicolumn{5}{@{}l}{\textit{Peldaño 4: multietapa (no implementado)}}\\
Árbol de modos & caudales por nodo del árbol del ACS, no anticipación & N2 o N3 & crece con el árbol & V8\\\bottomrule
\end{tabularx}
\caption{Las formulaciones de riesgo como modificaciones de N2. «Igual»: mismo tamaño y tiempo
que N2. Detalle de cada una en la sección indicada en la tabla~\ref{tab:riesgo-guia}.}
\label{tab:riesgo-formulaciones-n2}
\end{table}

\paragraph{Peldaño 1: el DSM.} Es una penalización sobre la PCT \emph{predicha nominal}. Con el
umbral sobre la consigna nunca se activa en N2 (sección~\ref{sec:dsm-lambda}); con el umbral bajo
la consigna, mueve la consigna. En ninguno de los dos casos sabe nada de la falla.

\paragraph{Peldaño 2, primera fuente: el pronóstico (F4).} La restricción gaussiana y la SAA de
dos etapas tratan la fuga futura como aleatoria. La SAA se montó sobre N3, con 59 trayectorias
del ensamble MC-dropout: el primer caudal es común y los demás se adaptan a cada trayectoria. Es
la formulación más grande de esta parte, y no cambió ninguna decisión.

\paragraph{Peldaño 2, segunda fuente: las fallas del ACS.} Las cuatro formulaciones de la
sección~\ref{sec:riesgo-formulaciones} usan la misma curva $\Delta T$ del ROM sin ACS. Las tres
primeras (cuantil, CVaR, Wilks) terminan en un techo $1220-m$; la cuarta (pérdida esperada) es
una SAA con 40 escenarios que no es un techo, pero cuya solución activa un margen efectivo. Las
fallas parciales dentro del horizonte son la única formulación en que la incertidumbre multiplica
a la decisión ($\eta u$), y por eso la única que podría no reducirse a un margen.

\paragraph{Peldaño 3: el escenario único.} Reemplaza la SAA de la pérdida esperada por un solo
escenario, elegido para que la decisión sea la misma. La versión aprendida usa un pronóstico
$\Phi(\theta,x)$ del margen, ajustado fuera de línea (figura~\ref{fig:f9-estrategia}).

\subsection{Consideraciones de cada formulación}
\begin{itemize}
  \item \textbf{Qué garantiza.} La cota con cuantil garantiza $P(\text{cruce}\mid\text{falla})\le\alpha$
        para una falla que empieza en la consulta, si la predicción corregida es exacta y valen A1--A2
        (sección~\ref{sec:prop-aditiva}); no garantiza nada sobre fallas que ya empezaron ni sobre
        varias fallas. Wilks garantiza lo mismo con $\alpha=0.05$ y confianza 0.952. El CVaR y la
        pérdida esperada no garantizan un nivel, sino una severidad o un equilibrio de costos.
  \item \textbf{Qué parámetro hay que elegir.} $\alpha$ es una probabilidad, comparable con los
        criterios del APS. $\beta$ y $n$ se traducen a un $\alpha$ equivalente (CVaR$_{0.95}\approx0.015$,
        Wilks $59\approx0.01$). $\rho$ no tiene una lectura física directa.
  \item \textbf{Qué supone sobre las fallas.} Todas suponen la distribución de la duración. Si las
        fallas reales duran más que las de diseño, la cota pierde su garantía (0.26 con
        $\alpha=0.02$); si duran menos, sobra pero no cuesta (sección~\ref{sec:riesgo-aplicabilidad}).
  \item \textbf{Qué supone sobre el modelo.} El calentamiento lo calcula el ROM. En la planta, los
        estados pre-enfriados se calientan algo más lento ($\sim10$\,\%), así que el error tiene el
        signo seguro (A1 conservador, sección~\ref{sec:prop-aditiva}).
  \item \textbf{Efectos de lazo cerrado.} Con la cota activa, el costo queda casi plano en la
        dirección del primer movimiento y el caudal oscila con período dos; $R\ge0.03$ lo corrige
        (sección~\ref{sec:supresion}). El escenario aprendido con el margen actual como contexto
        también oscila; con variables lentas, no (sección~\ref{sec:rozas-lazo}).
  \item \textbf{Cómputo.} Todas caben en los 0.5~s: las cotas en $\approx0.1$~s y la SAA de 40
        escenarios en $\le0.42$~s en el peor lazo.
\end{itemize}

\subsection{Por qué sobre N2 y cómo sería sobre N3}
La capa se montó sobre N2 por tres razones. (i) \emph{Da la misma decisión que N3}: en lazo
cerrado las trayectorias difieren en menos de 7~mK (sección~\ref{sec:trayectorias}), y la
decisión de N3 está certificada a menos del 0.5\,\% del óptimo global, con la misma garantía para
N2 (sección~\ref{sec:f7}); con $R=0.1$ siguen coincidiendo (menos de 0.002~K,
sección~\ref{sec:supresion}). (ii) \emph{Es convexo}: un QP admite holguras $\ell_1$, CVaR y SAA
sin perder el óptimo global, y el teorema de Homem-de-Mello exige convexidad. (iii) \emph{Es
rápido}: un QP de 77 variables en $\sim30$~ms deja tiempo para la simulación del ROM y para 40
escenarios.

Sobre N3, la cota sería una restricción no lineal más, $\hat y_k(x_k)+d\le1220-m$, que Ipopt
acepta sin cambios; la curva $\Delta T$ es la misma, porque se calcula con el ROM fuera del
optimizador. La SAA con N3 ya se hizo en F4 (2~s con 59 escenarios). Lo que se perdería es la
garantía de óptimo global de cada QP y la aplicabilidad directa del teorema del escenario único,
que necesita una segunda etapa lineal. Como la decisión de N2 y N3 coincide, los resultados de
riesgo valen para las dos.

\takeaway{\textbf{Resumen.} Todas las capas de riesgo dejan intacto el modelo de predicción y
agregan al QP de N2 un techo sobre la PCT (cuantil, CVaR, Wilks), un término de costo (DSM,
pérdida esperada) o un escenario (Homem-de-Mello). Lo que distingue a una métrica
probabilística de una determinista no es la forma del término, sino de dónde sale su constante:
del proceso de fallas o de un número fijo.}
